% Generated draft: requires mathematical and source review.
\documentclass[11pt]{article}
\usepackage[margin=25mm]{geometry}
\usepackage{fontspec}
\setmainfont{Latin Modern Roman}
\usepackage{amsmath,amssymb,mathtools,unicode-math}
\setmathfont{Latin Modern Math}
\usepackage{enumitem,needspace,xurl,hyperref}
\setlength{\parindent}{0pt}
\setlength{\parskip}{5pt}
\setlength{\emergencystretch}{3em}
\newcommand{\N}{\mathbb{N}}
\newcommand{\Z}{\mathbb{Z}}
\newcommand{\Q}{\mathbb{Q}}
\newcommand{\R}{\mathbb{R}}
\newcommand{\C}{\mathbb{C}}
\newcommand{\F}{\mathbb{F}}
\newcommand{\A}{\mathbb{A}}
\newcommand{\PP}{\mathbb{P}}
\newcommand{\E}{\mathbb{E}}
\newcommand{\eps}{\varepsilon}
\newcommand{\dd}{\,\mathrm{d}}
\begin{document}
\section*{UnsolvedMath ID: 11000042}
\textbf{Source problem code:} catalogue entry
\subsection{Definitions and mathematical statement}
Let $\mathcal M_g$ denote the genus-$g$ moduli space in the source's convention, and let $p:Y\to\mathcal M_g$ be any finite cover. A finite-order homeomorphism is a homeomorphism $f:Y\to Y$ such that $f^m=\operatorname{id}_Y$ for some integer $m\geq1$.

Question 4.3 asks whether, for every such cover and homeomorphism,
\[
 \bigl(f^m=\operatorname{id}_Y\text{ for some }m\geq1
 \ \text{ and }\ f\simeq\operatorname{id}_Y\bigr)
 \quad\Longrightarrow\quad f=\operatorname{id}_Y.
\]
Here $f\simeq\operatorname{id}_Y$ means there is a homotopy of maps from $f$ to the identity; an isotopy through homeomorphisms is not assumed. The conclusion is equality as maps. Neither regularity of the finite cover nor smoothness or preservation of a metric by $f$ is imposed.

The supplied excerpt does not fix the genus range or specify ordinary covers of the coarse moduli space versus orbifold covers of the moduli orbifold. In the latter convention, torsion-free finite-index subgroups of the mapping class group yield manifold covers of Teichm\"uller space modulo that subgroup, but the question does not explicitly restrict to them. These choices and potential low-genus exceptions must be resolved before asserting the quantified statement in one category.
\subsection{Short English statement}
Must every finite-order homeomorphism of a finite moduli-space cover that is homotopic to the identity actually equal the identity? The genus and covering category are unspecified.
\subsection{Sources}
\begin{itemize}
\item Local registry: \url{data/unsolvedmath/problems.json}, permanent UnsolvedMath ID $11000042$, exact source code \texttt{catalogue entry}.
\item Farb, \emph{Problems on Mapping Class Groups and Related Topics}, Question 4.3, PDF page 38: \url{https://www.math.uchicago.edu/~farb/papers/mcgbook.pdf}. Supplied access date: 2026-07-29; extraction: pdf-font-structure; page-disambiguated numbering among 329 headings.
\item No individual attribution or additional primary citation is supplied. The registry's literature assessment says ``Truncated. Cannot verify precisely.''
\end{itemize}
\end{document}
